% This supplementary chapter preserves the logarithmic-degeneration outlook.
% It is not part of the default fourteen-talk route.

\chapter{Degenerating domains and the virtual-class boundary}
\label[chapter]{chap:Degenerating_Domains_Virtual_Class_Boundary}

The preceding chapters reached a precise endpoint. For an elliptic equation
on a proper family of smooth compact domains, the intrinsic solution functor
is represented by a derived smooth object. Its tangent complex at a solution
is the deformation-obstruction complex of the linearized elliptic operator.
Steffens reached this conclusion through local Kuranishi charts, while Pardon
reorganized the proof around ordinary regularity and the extension of proper
section stacks from smooth to derived test objects.

This endpoint is substantial, but it is not yet the endpoint required by most
enumerative applications. A sequence of pseudo-holomorphic curves may develop
a long neck, bubble, or break. Its domains then leave the class of smooth
compact Riemann surfaces used in
\Cref{chap:Pseudo_Holomorphic_Curves_Derived_Moduli}.
The first question of this chapter is therefore geometric:

\begin{quote}
  How should a smooth family remember a domain whose neck becomes infinitely
  long?
\end{quote}

The elementary answer begins with the multiplication map
\[
  \R'_{\geq0}\times\R'_{\geq0}
  \longrightarrow
  \R'_{\geq0},
  \qquad
  (x,y)\longmapsto xy.
\]
Here $\R'_{\geq0}$ is the \emph{log ray}. Its underlying topological space is
the ordinary half-line, but the endpoint is treated as a point at cylindrical
infinity. If
\[
  x=e^{-s},\qquad y=e^{-t},\qquad \lambda=e^{-L},
\]
then the fiber $xy=\lambda$ is described by
\[
  s+t=L.
\]
Thus the positive fiber is an interval of length $L=-\log\lambda$, and the
limit $\lambda\to0$ is the limit in which this neck becomes infinitely long.
This one calculation will lead us to log smooth manifolds, exact
submersions, and log derived smooth manifolds.

The source status is part of the mathematics. Pardon's published 2024 article
develops the logarithmic examples and phrases the representability statement
as an expectation
\cite[Section~6]{Pardon_Representability_Elliptic_Fredholm}. His July 2026
manuscript states a Log Regularity Theorem and a Log Derived Regularity
Theorem, but explicitly describes itself as unfinished work in progress
\cite[P.6--P.7]{Pardon_Logarithmic_Derived_Moduli}. Its proposed proofs contain
unresolved references, and one proposition needed for the derived-base
reduction has no statement or proof. We will therefore separate three kinds
of assertion:

\begin{enumerate}
  \item Calculations and structural logarithmic theorems which are established.
  \item The analytic architecture of the gluing argument.
  \item Work-in-progress representability statements which organize the
    research frontier but will not be treated as completed results.
\end{enumerate}

The final question is logical rather than technical. Even a represented log
derived moduli stack need not be compact or oriented, and representability
does not automatically produce a virtual fundamental class. The chapter ends
by locating these further steps and by reviewing the path which led from
$\Cinfty$-rings to elliptic moduli problems.

\section{Why smooth domains are not the endpoint}
\label[section]{sec:Degeneration_Smooth_Domain_Boundary}

\subsection{What has already been represented}
\label[subsection]{sec:Degeneration_Represented_Endpoint}

Let
\[
  W\longrightarrow C\longrightarrow B
\]
be an elliptic section problem over a proper family of smooth compact domains.
In the setting of the last three chapters, its intrinsic solution stack is a
pullback
\[
\begin{tikzcd}
  \operatorname{Hol}_B(C,W) \rar \dar \drar[pullback]
    & \operatorname{Sec}_B(C,W) \dar["\Phi"] \\
  B \rar["0"'] & \operatorname{Sec}_B(C,H).
\end{tikzcd}
\]
The representability theorem identifies this functor locally with a derived
manifold. At a solution $(b,u)$, its relative tangent complex is represented
by the linearized elliptic operator
\[
  \left[
    \Gamma(C_b,u^*T_{W/C})
    \xrightarrow{D_{(b,u)}\Phi}
    \Gamma(C_b,u^*H)
  \right].
\]
The first term describes infinitesimal deformations and the second describes
equations. The kernel and cokernel give the classical deformation and
obstruction spaces.

Nothing in this statement adds new objects to the family of domains. It says
that the solution functor for the specified family is represented, including
at nonregular solutions. It does not say that a sequence of solutions remains
inside this functor when its domains degenerate.

This distinction can be phrased abstractly. Derived geometry improves the
intersection theory of objects which are already present. Compactification
enlarges the collection of objects. If a moduli functor contains only smooth
compact domains, no amount of derived structure on that functor can turn a
nodal domain into one of its points.

\subsection{A neck which escapes}
\label[subsection]{sec:Degeneration_Escaping_Neck}

Consider a family of Riemann surfaces containing an embedded annular region.
As a parameter approaches a boundary value, the annulus becomes longer and
thinner. In conformal coordinates it looks like
\[
  [0,L]\times S^1
\]
with $L\to\infty$. The limiting picture consists of two surfaces with
half-infinite cylindrical ends. Their ends must be matched if they are to be
remembered as the two branches of one node.

For a sequence of pseudo-holomorphic maps, this geometric degeneration may
interact with the maps themselves. Energy can concentrate, new sphere
components can bubble off, and stability conditions may be needed to obtain a
finite automorphism group. These phenomena lead to different compactified
moduli problems in Gromov--Witten theory, symplectic field theory, and related
constructions.

We will not construct any one of these compactifications. The purpose of the
neck is more basic: it shows which differentiable structure a family of
degenerating domains ought to carry. The relevant structure is already
visible in a one-dimensional real model.

\begin{warning}[Category is not compactness]
\label[warning]{warn:Degeneration_Category_Not_Compactness}
A category capable of expressing nodal or broken objects does not prove that
a chosen moduli problem is compact in that category. One must still specify
the limiting objects and prove a compactness theorem.
\end{warning}

\section{One neck in complete detail}
\label[section]{sec:Degeneration_One_Neck}

\subsection{The multiplication family}
\label[subsection]{sec:Degeneration_Multiplication_Family}

Begin with the continuous map
\[
  \mu\colon\R_{\geq0}^2\longrightarrow\R_{\geq0},
  \qquad
  \mu(x,y)=xy,
\]
and restrict the source to the square $0\leq x,y\leq1$. For
$\lambda>0$, the fiber
\[
  C_\lambda=\{(x,y):xy=\lambda\}
\]
is a smooth arc. Its endpoints are $(1,\lambda)$ and $(\lambda,1)$. The
central fiber is
\[
  C_0=\{(x,y):xy=0\},
\]
the union of the two coordinate axes.

The topological picture is familiar, but it hides the metric length of the
neck. On the positive locus, introduce logarithmic coordinates
\begin{equation}
  \label{eq:Degeneration_Log_Coordinates}
  x=e^{-s},
  \qquad
  y=e^{-t},
  \qquad
  \lambda=e^{-L}.
\end{equation}
Then
\begin{equation}
  \label{eq:Degeneration_Multiplication_Addition}
  xy=\lambda
  \quad\Longleftrightarrow\quad
  s+t=L.
\end{equation}
The inequalities $0\leq x,y\leq1$ become $0\leq s,t\leq L$. Projection to
$s$ therefore identifies $C_\lambda$ with the interval $[0,L]$.

\begin{observation}[Multiplication becomes addition]
\label[observation]{obs:Degeneration_Multiplication_Addition}
The gluing parameter is $\lambda$, while the neck length is
\[
  L=-\log\lambda.
\]
Multiplication of boundary parameters becomes addition of cylindrical
lengths.
\end{observation}

As $\lambda\to0$, the length $L$ tends to infinity. From the left endpoint,
the coordinate $s$ runs along a half-infinite interval. From the right
endpoint, the coordinate $t$ does the same. The limiting geometric picture is
therefore a pair of cylindrical ends.

The basic central fiber $C_0$ is not literally a disjoint union: its two axes
meet at the origin. The log structure remembers the two branches, their
cylindrical directions, and the common degeneration parameter. In an enlarged
``punctured'' logarithmic category, one may subsequently split suitable unions
of strata, separating the two axes while retaining their matching at infinity
\cite[P.6.2--P.6.3]{Pardon_Logarithmic_Derived_Moduli}. The basic log fiber and
this separated version should not be conflated.

\subsection{The complex nodal model}
\label[subsection]{sec:Degeneration_Complex_Node}

The real multiplication family is the radial part of the standard complex
nodal family
\begin{equation}
  \label{eq:Degeneration_Complex_Node}
  zw=q.
\end{equation}
Write
\[
  z=xe^{i\theta},
  \qquad
  w=ye^{i\varphi},
  \qquad
  q=\lambda e^{i\psi}.
\]
Equation~\eqref{eq:Degeneration_Complex_Node} is equivalent to the two
equations
\begin{equation}
  \label{eq:Degeneration_Radial_Angular_Matching}
  xy=\lambda,
  \qquad
  \theta+\varphi=\psi.
\end{equation}
The first controls the neck length. The second controls the angular matching
of the two circles at infinity.

This is the oriented real blow-up of the complex nodal family. Instead of
collapsing all angular directions at $z=w=0$, it records the circle along
which the two cylindrical boundaries are identified. For the purposes of this
chapter, this is all that is needed from the geometry of stable curves. The
formula explains why a degeneration contains both a length parameter and
matching data.

\subsection{Why ordinary differentiability is not intrinsic}
\label[subsection]{sec:Degeneration_Log_Differentiability_Motivation}

The coordinate change in \Cref{eq:Degeneration_Log_Coordinates} satisfies
\[
  \frac{dL}{d\lambda}=-\frac1\lambda.
\]
Thus a bounded variation of the cylindrical length is singular when measured
by the ordinary vector field $\partial_\lambda$ near $\lambda=0$. The natural
relation is instead
\begin{equation}
  \label{eq:Degeneration_Log_Vector_Field}
  \lambda\partial_\lambda=-\partial_L.
\end{equation}
Translation along a long cylinder is controlled by $\partial_L$, so the
vector field visible to the geometry is $\lambda\partial_\lambda$.

This is not a claim that the underlying topology of a manifold with boundary
is inadequate. The ordinary half-line and the log ray have the same
underlying topological space. The distinction lies in the tangent theory and
in which parameter-dependent functions are declared smooth. Gluing maps which
are natural in the length $L$ need not be ordinarily smooth in
$\lambda=e^{-L}$, while they can be logarithmically smooth.

The phenomenon is familiar in analysis on manifolds with cylindrical ends.
Solutions and their derivatives often converge as $L\to\infty$, and the error
terms decay exponentially in $L$. A differentiability theory built from
$\partial_L$ records precisely this behavior.

\section{The minimum logarithmic language}
\label[section]{sec:Degeneration_Minimum_Log_Language}

\subsection{The log ray and asymptotic smoothness}
\label[subsection]{sec:Degeneration_Log_Ray}

\begin{definition}[Log ray]
\label[definition]{def:Degeneration_Log_Ray}
The \emph{log ray} $\R'_{\geq0}$ has underlying topological space
$\R_{\geq0}$. On its positive locus, the coordinate $x=e^{-s}$ identifies it
with a cylinder coordinate $s\in\R$. The endpoint $x=0$ is regarded as the
asymptotic point $s=\infty$, and its logarithmic tangent and cotangent
directions are generated by
\[
  x\partial_x=-\partial_s,
  \qquad
  \frac{dx}{x}=-ds.
\]
\end{definition}

One way to recognize the resulting smooth structure is through asymptotics.
A log smooth function
\[
  f\colon\R^n\times\R'_{\geq0}\longrightarrow\R
\]
has, in cylindrical coordinates, the form
\[
  f(u,s)=f_\infty(u)+o(1)_{C^\infty}
  \qquad (s\to\infty).
\]
Here all derivatives in $u$ and $s$ of the error converge to zero uniformly
over compact subsets of the ordinary variables. A log smooth map to the log
ray has logarithmic coordinate of the form
\[
  g(u,s)=g_\infty(u)+a s+o(1)_{C^\infty}
\]
for a constant $a\geq0$. The linear term records the monomial behavior in the
logarithmic direction
\cite[Definition~6.1 and Example~6.2]
  {Pardon_Representability_Elliptic_Fredholm}.

For elliptic gluing, one often needs the stronger estimate
\[
  O(e^{-\delta s})_{C^\infty}
\]
for some $\delta>0$. This distinction is essential. Log smoothness describes
the ambient geometric category. Exponential decay is an analytic conclusion
about solutions, or additional structure imposed on a family. It is not part
of the bare definition of a log smooth manifold.

\subsection{Toric charts, strata, and depth}
\label[subsection]{sec:Degeneration_Toric_Charts}

The log ray is the simplest member of a family of real toric models. Let $P$
be a real polyhedral cone and define
\begin{equation}
  \label{eq:Degeneration_Real_Toric_Model}
  X_P=
  \Hom\bigl((P,+),(\R_{\geq0},\cdot)\bigr).
\end{equation}
The faces of $P$ stratify $X_P$. For $P=\R_{\geq0}$, one obtains
$X_P=\R'_{\geq0}$.

\begin{definition}[Log smooth manifold]
\label[definition]{def:Degeneration_Log_Smooth_Manifold}
A \emph{log smooth manifold} is a log topological space admitting an atlas by
open subsets of the $X_P$, with transition maps which are log smooth. Locally,
such a map is monomial in the logarithmic directions and smooth in the
ordinary directions.
\end{definition}

The precise definition of the underlying log structure is given in the
supplement below. For the live route, the important datum at a point
$m\in M$ is its characteristic cone $\mathcal Z_{M,m}$. It records which
stratum contains $m$ and which logarithmic directions leave that stratum. Its
dimension is the \emph{depth} of $m$.

Depth-zero points form an ordinary smooth manifold. A depth-one point may be
viewed as a point at infinity of a cylindrical end. A map can also have
relative depth one, as in the multiplication family: the total space acquires
one more logarithmic direction than the base at the central fiber.

Topologically, an orthant chart is a manifold with corners. Its logarithmic
tangent theory is nevertheless different. At a boundary point of the
ordinary half-line, the vector $x\partial_x$ vanishes as an ordinary vector
field but generates the logarithmic tangent direction. This is why log smooth
manifolds are not merely manifolds with corners under a new name.

\subsection{Strict and exact submersions}
\label[subsection]{sec:Degeneration_Strict_Exact_Submersions}

A morphism of log smooth manifolds induces a map on characteristic cones. It
is \emph{strict} if this map is an isomorphism at every point. Thus a strict
map changes only the ordinary smooth directions and preserves the local
logarithmic combinatorics. A strict submersion is locally an ordinary smooth
projection over a fixed toric chart.

Exact maps allow the characteristic cones to change, but in a controlled way.
The formal cone-theoretic definition appears in the supplement. Its operative
meaning here is that logarithmic fibers have the expected strata and remain
well behaved under base change.

The local models explain the terminology. An exact submersion of relative
depth at most one is locally built from pullbacks of
\begin{equation}
  \label{eq:Degeneration_Depth_One_Local_Models}
  \R^k\longrightarrow *,
  \qquad
  \R'_{\geq0}\longrightarrow *,
  \qquad
  (\R'_{\geq0})^2
    \xrightarrow{\,\mu\,}
  \R'_{\geq0}.
\end{equation}
The first is an ordinary smooth parameter. The second is a cylindrical end.
The third is the neck degeneration of \Cref{sec:Degeneration_One_Neck}.

\begin{theorem}[Pullback stability of exact submersions]
\label[theorem]{thm:Degeneration_Exact_Submersions_Pullback}
Exact submersions of log smooth manifolds are preserved under pullback.
\end{theorem}

This is Pardon's P.6.4, proved as T.9.6.17
\cite[P.6.4 and T.9.6.17]{Pardon_Logarithmic_Derived_Moduli}. In the running
example, a map
\[
  \lambda\colon Z\longrightarrow\R'_{\geq0}
\]
pulls the multiplication family back to
\[
  \{(z,x,y):xy=\lambda(z)\}
  \longrightarrow Z.
\]
The theorem says that this is again an exact degenerating family, including
over points at which $\lambda(z)=0$.

The general proof is local on source and target. A submersion is locally
monomial in the logarithmic variables, so the pullback reduces to a pullback
of real toric charts. Exactness ensures that the cone-theoretic fiber product
has the expected underlying monoid and strata. Ordinary smooth directions are
then added by products. We will use the theorem, rather than reproduce its
full cone calculation.

\subsection{Deriving only the smooth directions}
\label[subsection]{sec:Degeneration_Log_Derived_Smooth_Manifolds}

One could formally derive the whole category of log smooth manifolds. Pardon
denotes this unrestricted derived completion by $D(\operatorname{LogSm})$.
It is not the category selected by the elliptic moduli problem.

\begin{definition}[Log derived smooth manifold]
\label[definition]{def:Degeneration_Log_Derived_Smooth_Manifold}
The $\infty$-category $\operatorname{LogDSm}$ is the full subcategory of
$D(\operatorname{LogSm})$ whose objects are locally finite limits
\[
  \lim_{k\in K}p(k)
\]
of finite diagrams $p\colon K^{\triangleright}\to\operatorname{LogSm}$ for
which every map
\[
  p(k)\longrightarrow p(*)
\]
is a strict submersion.
\end{definition}

The definition derives the smooth equations while leaving the logarithmic
combinatorics geometric. In particular, every log derived smooth manifold is
locally strict over some toric chart $X_P$. Its smooth directions may carry a
nontrivial tangent complex, while the characteristic cone still records an
ordinary stratified degeneration parameter.

This restricted definition has the preservation property required by the
neck family.

\begin{theorem}[Exact pullbacks in log derived smooth geometry]
\label[theorem]{thm:Degeneration_LogDSm_Exact_Pullbacks}
The inclusion
\[
  \operatorname{LogSm}\longrightarrow\operatorname{LogDSm}
\]
preserves exact submersive pullbacks.
\end{theorem}

This is Pardon's P.6.6, proved as T.10.2.4
\cite[P.6.5--P.6.6 and T.10.2.4]
  {Pardon_Logarithmic_Derived_Moduli}. It is the logarithmic analogue of the
transverse-pullback principle which characterized derived manifolds in
\Cref{chap:From_Smooth_Algebra_To_Derived_Manifolds}. The category is chosen so
that the pullbacks intrinsic to the geometry remain the correct pullbacks
after passage to derived objects.

\section{The proposed logarithmic regularity route}
\label[section]{sec:Degeneration_Proposed_Log_Regularity}

\subsection{Elliptic equations over a degenerating family}
\label[subsection]{sec:Degeneration_Log_Elliptic_Problem}

Consider a diagram
\begin{equation}
  \label{eq:Degeneration_Log_Elliptic_Family}
  W\longrightarrow C\longrightarrow B
\end{equation}
in log smooth manifolds. Suppose that $W\to C$ is strict and that $C\to B$
is a proper exact submersion of relative depth one. The standard locus of a
fiber is an ordinary, possibly noncompact, smooth manifold with cylindrical
ends. Its logarithmic bordification is compact because its points at infinity
have been included.

The solution stack is defined by the same diagrammatic prescription as
before. On a test object $T\to B$, a point consists of a section of
$C\times_BT$ into $W$ satisfying the vertical differential equation. The
definition makes sense on log topological, log smooth, and log derived smooth
test objects. What changes is the analysis near the relative depth-one locus.

The differential equation has an asymptotic limit on the cylindrical ends.
For a family containing nodes, this asymptotic problem contains both the ends
which persist in each fiber and the two branches created by a degeneration.
A solution is \emph{regular at infinity} when this asymptotic problem is
regular relative to the base. Full regularity additionally requires the
linearized operator on the whole domain to be surjective.

In the nondegenerate case, the asymptotic operator has no unwanted kernel. In
a Morse--Bott situation, the asymptotic solutions form a smooth family, and
one needs a fiberwise exponential-decay structure to control motion along it.
These distinctions are important in the source, but the nondegenerate case is
enough to display the argument.

\subsection{What a gluing theorem must accomplish}
\label[subsection]{sec:Degeneration_Gluing_Architecture}

Suppose first that the degeneration parameter is zero. A solution on the
broken fiber is represented by solutions $u_-$ and $u_+$ on two cylindrical
pieces, with matching asymptotic values. For large finite $L$, truncate the
two ends and splice them along a cylinder of length $L$. The result is a
preglued section $u_L^{\mathrm{pre}}$.

Pregluing is not expected to solve the nonlinear equation exactly. Its defect
is concentrated in the cutoff region. Exponential convergence on the ends
gives a schematic estimate
\begin{equation}
  \label{eq:Degeneration_Pregluing_Error}
  \|\Phi(u_L^{\mathrm{pre}})\|
  \leq C e^{-\delta L}.
\end{equation}
To correct the approximate solution, one studies the linearization at
$u_L^{\mathrm{pre}}$. A gluing theorem must supply four ingredients.

\begin{enumerate}
  \item Exponential decay of the original solutions and all derivatives on
    their cylindrical ends.
  \item Uniform control of the family of linearized operators as
    $L\to\infty$.
  \item A right inverse with norm bounded independently of sufficiently large
    $L$, after the appropriate matching conditions are imposed.
  \item A Newton--Picard argument which corrects $u_L^{\mathrm{pre}}$ to a
    unique nearby exact solution.
\end{enumerate}

The inverse construction cuts an exact solution on a long neck into its two
limiting pieces. Together, the two directions identify a neighborhood of the
broken configuration with a neighborhood in the gluing parameter and the
moduli of limiting solutions.

The logarithmic content is stronger than continuity. The exact solutions must
depend log smoothly on the parameter $\lambda=e^{-L}$, including in families.
Thus the estimates must control derivatives in ordinary base variables and in
the logarithmic derivative $\lambda\partial_\lambda=-\partial_L$.

We will not prove these analytic claims. They are the additional input which
distinguishes ordinary regularity on a fixed compact domain from regularity
across a cylindrical degeneration. Pardon's 2024 article points to the gluing
analysis of Fukaya--Oh--Ohta--Ono as the essential precedent
\cite[Section~6]{Pardon_Representability_Elliptic_Fredholm}. The 2026
manuscript proposes a general treatment in N.5.3
\cite[N.5.3]{Pardon_Logarithmic_Derived_Moduli}.

\subsection{The announced regularity statements}
\label[subsection]{sec:Degeneration_Announced_Regularity}

The following discussions record statements in the July 2026 manuscript.
Their status is included in their titles.

\begin{discussion}[Work-in-progress Log Regularity Theorem]
\label[discussion]{disc:Degeneration_Log_Regularity_WIP}
For a proper exact depth-one family of elliptic section problems over a log
smooth base, the regular locus of the log smooth solution stack is represented
by a log smooth manifold, and its comparison with the corresponding log
topological solution stack is an isomorphism. A Morse--Bott variant is stated
after choosing suitable fiberwise exponential-decay data.
\end{discussion}

This is P.7.2, with its proposed proof in N.5.3.1
\cite[P.7.2 and N.5.3.1]{Pardon_Logarithmic_Derived_Moduli}. The new analytic
content is precisely the gluing architecture of the preceding subsection.

\begin{discussion}[Work-in-progress Log Derived Regularity Theorem]
\label[discussion]{disc:Degeneration_Log_Derived_Regularity_WIP}
Under the corresponding nondegenerate or Morse--Bott hypotheses at infinity,
the regular-at-infinity locus of the log derived solution stack is
representable, and its comparison with the log topological solution stack is
an isomorphism.
\end{discussion}

This is the content of P.7.3, elaborated as N.6.5.3
\cite[P.7.3 and N.6.5.3]{Pardon_Logarithmic_Derived_Moduli}. The distinction
between the two statements is important. The first assumes regularity of the
full elliptic problem and concludes an ordinary log smooth moduli manifold.
The second requires regularity only at infinity. Interior cokernel directions
are retained as derived obstruction data.

Neither discussion is a theorem proved in these notes. They describe the
intended endpoint and let us understand why the structural logarithmic
theorems were designed as they were.

\subsection{The route from regularity to derived representability}
\label[subsection]{sec:Degeneration_Log_Derived_Route}

The proposed proof mirrors \Cref{chap:Pardon_Derived_Regularity}. Its stages
can be aligned as follows.

\begin{enumerate}
  \item Ordinary parametric regularity is replaced by log regularity, whose
    new input is gluing across the depth-one locus.
  \item Smooth section stacks are replaced by log smooth section stacks which
    retain the asymptotic data of the family.
  \item Pullbacks along submersions are replaced by strict or exact
    submersive pullbacks, according to whether the logarithmic combinatorics
    remain fixed.
  \item Finite-dimensional parameter thickening is performed exactly as in
    the smooth case, but the perturbations are supported away from the ends.
  \item The original problem is recovered as a log derived zero-parameter
    fiber of the regular enlarged problem.
\end{enumerate}

Let us isolate the fourth step. Fix a solution $(b,u)$. Its interior
linearization is Fredholm. Choose a finite-dimensional vector space $E$ and a
map
\[
  \alpha\colon E\longrightarrow\Gamma(C_b,u^*H)
\]
whose image surjects onto the relevant cokernel. The representatives can be
chosen with compact support in the open locus of $C$ where the family is
strict over $B$. Extend them to nearby sections and define
\[
  \Phi^+(u,e)=\Phi(u)+\alpha(e).
\]
The enlarged equation is regular at $(b,u,0)$ in the interior. Because
$\alpha$ vanishes near the cylindrical ends, it does not change the
asymptotic operator, the matching condition, or regularity at infinity.

If the enlarged regular locus is represented by a log smooth manifold
$M^+_{\mathrm{reg}}$ over $B^+=B\times E$, the original local moduli problem
is recovered as
\begin{equation}
  \label{eq:Degeneration_Log_Derived_Zero_Parameter_Fiber}
  M^+_{\mathrm{reg}}
  \mathbin{\times^R}_{B^+}B.
\end{equation}
The smooth equations in this fiber may be nontransverse, so the pullback is
derived. The map $B\to B^+$ is strict in the logarithmic directions, so the
underlying degeneration parameter is unchanged. This illustrates the slogan:

\begin{quote}
  Derive the smooth equation; retain the logarithmic degeneration.
\end{quote}

N.6.5.1 contains this parameter-thickening step and observes explicitly that
the perturbations should have compact support in the strict locus
\cite[N.6.5.1]{Pardon_Logarithmic_Derived_Moduli}.

\subsection{Where the proposed proof is incomplete}
\label[subsection]{sec:Degeneration_Log_Source_Status}

The architecture is persuasive, but the July 2026 source does not close its
dependencies. The central points are these.

\begin{enumerate}
  \item The title page describes the manuscript as an excerpt of unfinished
    work in progress.
  \item N.5.3.1 contains a substantial proposed proof of log smooth
    regularity, but invokes unresolved comparisons and analytic assertions.
  \item N.6.5.1 proves the parameter-thickening cover in the logarithmic
    setting.
  \item N.6.5.2 consists only of the heading ``Proposition (compare
    N.6.3.1)'' followed by ``Proof.'' It contains neither a proposition nor a
    proof.
  \item N.6.5.3 invokes N.6.5.2 to reduce the problem over a log derived base
    to the case of a log smooth base.
  \item The remainder of N.6.5.3 also invokes unresolved references for a
    jet-level presentation of the equation, the extension of section stacks,
    sharp local presentations, and analytic-geometric tangent comparisons.
\end{enumerate}

Consequently, \Cref{disc:Degeneration_Log_Regularity_WIP,disc:Degeneration_Log_Derived_Regularity_WIP}
should be described as work-in-progress theorems. The manuscript supplies a
detailed program and proves important structural pieces, but it is not yet a
source in which the logarithmic representability argument can be read as a
closed proof.

This conclusion does not diminish the value of the framework. It determines
what can responsibly be taught from it. The neck model, the logarithmic
category, exact pullback stability, and the proposed proof architecture form
a coherent mathematical story. The final representability claim remains part
of the frontier.

\section{Compactness, virtual classes, and the seminar boundary}
\label[section]{sec:Degeneration_Virtual_Class_Boundary}

\subsection{Representability is not compactification}
\label[subsection]{sec:Degeneration_Representability_Not_Compactification}

The phrase ``logarithmic compactification'' can conceal three different
tasks. They should be separated.

\begin{enumerate}
  \item \emph{Choose an enlarged moduli problem.} Specify which nodal, broken,
    punctured, bubbled, or expanded objects are admitted, together with their
    morphisms, equivalences, and stability conditions.
  \item \emph{Prove compactness.} Show that every sequence satisfying the
    relevant bounds has a convergent subsequence in the enlarged topology.
  \item \emph{Prove representability.} Show that families of the chosen
    objects are represented in the desired smooth, logarithmic, or derived
    geometric category.
\end{enumerate}

Logarithmic geometry contributes to the first and third tasks. It supplies a
natural language for cylindrical ends, nodes, matching data, and expanded
targets. A logarithmic regularity theorem would show that the resulting
families have the expected differentiable structure near the boundary.

Compactness is a separate analytic and topological theorem. For
pseudo-holomorphic curves it uses energy bounds, removal of singularities,
bubbling analysis, and stability. The properness assumption on
$C\to B$ in the regularity theorem says that each logarithmically bordified
domain is compact over the parameter. It does not say that the solution stack
is proper.

Thus \Cref{warn:Degeneration_Category_Not_Compactness} remains in force even
if the work-in-progress representability theorem is completed.

\subsection{From a moduli functor to an invariant}
\label[subsection]{sec:Degeneration_Moduli_To_Invariant}

The full dependency chain is
\begin{equation}
  \label{eq:Degeneration_Virtual_Class_Ledger}
  \begin{gathered}
    \text{moduli functor}
    \longrightarrow
    \text{derived representability}
    \longrightarrow
    \text{compactification}
    \\
    \longrightarrow
    \text{orientation and virtual class}
    \longrightarrow
    \text{invariant}.
  \end{gathered}
\end{equation}
Each arrow requires new input.

The moduli functor specifies the objects, families, and equivalences. Derived
representability supplies a coherent geometric object and its tangent
complex. If this complex is two-term, it gives a virtual dimension
\[
  \operatorname{vdim}_{(b,u)}
  =\dim\ker D_{(b,u)}\Phi-dim\operatorname{coker}D_{(b,u)}\Phi.
\]
It also determines the determinant line
\[
  \det(D_{(b,u)}\Phi)
  =
  \Lambda^{\mathrm{top}}\ker D_{(b,u)}\Phi
  \otimes
  \bigl(\Lambda^{\mathrm{top}}
    \operatorname{coker}D_{(b,u)}\Phi\bigr)^*.
\]
Neither datum provides an orientation automatically.

Compactification adds the limiting objects needed for properness. An
orientation chooses a coherent orientation of the determinant data. A
virtual-class formalism then converts the compact oriented derived or
Kuranishi-type structure into a homology, cohomology, or bordism class of the
virtual dimension. Finally, evaluation maps and cohomological insertions turn
that class into the desired invariant.

There are several possible formalisms for the virtual-class step. Their
hypotheses and targets differ, so there is no responsible theorem saying that
an arbitrary represented derived smooth stack automatically has a virtual
fundamental class. The precise category, finiteness conditions, properness,
orientation, and chosen homology or bordism theory all matter.

\subsection{The comparison with implicit atlases}
\label[subsection]{sec:Degeneration_Implicit_Atlas_Comparison}

Pardon's implicit-atlas theory addresses the virtual-class problem by a route
which is logically separate from his derived representability program. An
implicit atlas is, roughly, a coherent system of finite-dimensional reductions
on a topological moduli space. An algebraic and sheaf-theoretic package turns
this data into virtual cochain complexes and virtual fundamental classes
\cite[Sections~1.1--1.3]{Pardon_Virtual_Fundamental_Cycles}.

The analytic input is correspondingly different. The implicit-atlas
construction requires suitable topological gluing results but does not require
the transition maps between gluing charts to be smooth. Logarithmic
representability seeks a stronger conclusion: the family itself should be
represented by a log smooth or log derived smooth object.

The outputs are also different.

\begin{enumerate}
  \item Derived representability produces an intrinsic geometric object with
    coherent deformation-obstruction theory.
  \item The implicit-atlas package produces virtual homological data from a
    compactified topological moduli problem equipped with finite-dimensional
    reductions.
\end{enumerate}

The shared Fredholm reductions do not make one construction a reformulation
of the other. A future comparison theorem could relate their outputs under
appropriate hypotheses, but no such comparison is needed for the narrative of
this seminar.

\section{Summary and supplementary materials}
\label[section]{sec:Degeneration_Summary_Supplementary_Materials}

\subsection{The arc of the seminar}
\label[subsection]{sec:Degeneration_Seminar_Arc}

We can now see the route from the first examples to the present frontier.

\begin{enumerate}
  \item $\Cinfty$-rings encode smooth operations and smooth equations
    algebraically.
  \item Derived pullbacks retain the infinitesimal information lost by an
    ordinary nontransverse intersection.
  \item The universal property of derived manifolds organizes these
    pullbacks coherently and without choosing a particular presentation.
  \item Stacks encode families, descent, and automorphisms.
  \item Relative jet bundles formulate differential equations intrinsically
    on families of sections.
  \item Ellipticity turns the infinite-dimensional linearization into finite
    deformation-obstruction data.
  \item Fredholm reduction and representability identify the intrinsic
    solution functor with a derived smooth moduli object.
  \item Degenerating domains require logarithmic geometry and gluing, while
    compactness, orientations, and virtual classes remain additional inputs.
\end{enumerate}

The same geometric operation has appeared at several scales. We formed
pullbacks to intersect finite-dimensional manifolds, to take zero loci of
sections, to define solution stacks of differential equations, and to change
parameters in families of cylindrical degenerations. The ambient category was
chosen each time so that the pullbacks intrinsic to the moduli problem were
the pullbacks it preserved.

This perspective also explains the homotopy-theoretic character of the
seminar. $\infty$-categories were not introduced as a separate formal subject.
They supplied the language in which universal properties, descent, derived
pullbacks, and representing objects can be stated without losing their
coherence.

\subsection{The live route}
\label[subsection]{sec:Degeneration_Live_Route}

The 90-minute talk can be compressed into the following steps.

\begin{enumerate}
  \item Smooth-domain representability does not add nodal or broken domains.
  \item The family $xy=\lambda$ models a neck degeneration.
  \item In logarithmic coordinates it becomes $s+t=L$ with
    $L=-\log\lambda$.
  \item The complex equation $zw=q$ adds angular matching at infinity.
  \item The log ray treats its endpoint as a point at cylindrical infinity.
  \item Log smooth manifolds are locally modelled on real toric charts.
  \item Exact depth-one submersions are locally built from ordinary
    projections, cylindrical ends, and multiplication necks.
  \item Exact submersions are stable under pullback.
  \item Log derived smooth geometry derives smooth equations while retaining
    the logarithmic degeneration.
  \item Gluing requires exponential decay, uniform inversion, and nonlinear
    correction on long necks.
  \item Pardon's logarithmic representability theorem is presently a
    work-in-progress claim.
  \item Logarithmic parameter thickening reuses the Fredholm argument of
    \Cref{chap:Pardon_Derived_Regularity} away from the ends.
  \item Logarithmic representability does not prove compactness.
  \item A virtual class additionally requires a compactified problem,
    orientation data, and a chosen virtual-class or bordism formalism.
\end{enumerate}

The final board should retain exactly three statements.

\begin{enumerate}
  \item The multiplication map of log rays turns a vanishing gluing parameter
    into a neck of diverging length.
  \item Log derived smooth geometry preserves the logarithmic pullbacks needed
    for degenerating families while retaining derived smooth obstruction data.
  \item Compactification and virtual classes require additional input, and
    the proposed logarithmic regularity theorem remains unfinished in the
    July 2026 manuscript.
\end{enumerate}

\subsection{Supplement: log structures and exactness}
\label[subsection]{sec:Degeneration_Supplement_Log_Structures}

For completeness, we record the formal layer suppressed in the live route. A
log structure on a topological space $X$ consists of a sheaf of commutative
monoids $\mathcal O_X^{\geq0}$, a morphism
\[
  \mathcal O_X^{\geq0}
  \longrightarrow
  \mathcal C_X^{\geq0}
\]
to the sheaf of nonnegative continuous functions under multiplication, and a
condition saying that this map is an isomorphism over the subsheaf of
everywhere positive functions. A morphism of log spaces is a continuous map
together with a compatible pullback map of these sheaves.

The toric space $X_P$ of \Cref{eq:Degeneration_Real_Toric_Model} carries the
log structure associated to the tautological map
\[
  P\longrightarrow\mathcal C_{X_P}^{\geq0}.
\]
After quotienting the stalk of the log structure by its invertible elements,
one obtains the characteristic monoid or, in Pardon's real setting, the
characteristic cone.

For a map $f\colon Q\to P$ of polyhedral cones, exactness means
\[
  (f^{\mathrm{gp}})^{-1}(P)=Q
\]
inside $Q^{\mathrm{gp}}$. A log map is exact when its induced maps on
characteristic cones are exact. This condition is inherited under localization
at faces and is designed to make logarithmic fiber products have the expected
stratified topology.

The pullback theorem
\Cref{thm:Degeneration_Exact_Submersions_Pullback} is the payoff. Locally, an
exact submersion is reduced to a monomial map of toric charts. The cone
fiber-product calculation produces the logarithmic part of the pullback, and
the ordinary submersion theorem handles the smooth variables. At relative
depth at most one, this reduces to the three models in
\Cref{eq:Degeneration_Depth_One_Local_Models}.

\subsection{Supplement: gluing profiles}
\label[subsection]{sec:Degeneration_Supplement_Gluing_Profiles}

An alternative way to place an ordinary smooth structure on a gluing
parameter is to choose a \emph{gluing profile}, a homeomorphism
\[
  \rho\colon\R'_{\geq0}\longrightarrow\R_{\geq0}
\]
which converts the desired asymptotic functions into ordinarily smooth
functions. For exponentially decaying errors, one possible profile is
\[
  \rho(x)=(-\log x)^{-1}.
\]
In the length coordinate this is $L\mapsto L^{-1}$. Such a profile can turn a
particular class of log smooth charts into charts on manifolds with corners.

The choice is useful but not neutral. It changes tangent coordinates and may
respect tangent bundles only after passing to an appropriate homotopical
comparison. The logarithmic formulation keeps $L$ itself as the intrinsic
cylindrical coordinate and records monomial changes of degeneration
parameters directly. This is one reason it is well suited to simultaneous
source and target degenerations.

\subsection{Supplement: detailed source audit}
\label[subsection]{sec:Degeneration_Supplement_Source_Audit}

The published article and the July 2026 manuscript should be read together.
The former is stable but deliberately brief. Definition~6.1 and
Examples~6.2--6.7 introduce log smooth manifolds, cylindrical coordinates,
exponential decay, and the multiplication family. Its term ``simply broken
submersion'' corresponds to the principal depth-one degeneration considered
here. The article concludes that one could hope to prove representability on a
suitable $\infty$-category of derived log smooth manifolds
\cite[Section~6]{Pardon_Representability_Elliptic_Fredholm}.

The July 2026 manuscript replaces this preliminary terminology by exact
submersions and develops the categorical framework in P.6, T.9, and T.10.
The pullback theorems used in this chapter have proofs in T.9.6.17 and
T.10.2.4. These structural statements are therefore on firmer ground than the
regularity results.

P.7.2 and P.7.3 state the two regularity theorems. N.5.3.1 develops the
ordinary log smooth argument using local linear models, finite-dimensional
projections, inverse operators, and Newton iteration. The proof still invokes
unresolved claims, including an identification of regularity in a local
splitting and smooth dependence of the inverse family.

N.6.5 is shorter. N.6.5.1 adapts parameter thickening by choosing perturbations
supported away from the depth-one locus. N.6.5.2 is empty. N.6.5.3 then cites
N.6.5.2 in its first reduction and continues through several additional
placeholders. The result is a blueprint, not a completed proof. Any later
version of the manuscript should be checked afresh before these status labels
are revised.

\subsection{Related literature}
\label[subsection]{sec:Degeneration_Related_Literature}

For the live route, the most efficient source is Section~6 of Pardon's
\emph{Representability in non-linear elliptic Fredholm analysis}. It contains
the log ray, the asymptotic description of log smooth maps, punctured Riemann
surfaces, exponential decay, the oriented real blow-up of $zw=q$, and the
representability expectation
\cite[Section~6]{Pardon_Representability_Elliptic_Fredholm}.

P.6--P.7 of
\emph{Logarithmic derived moduli theory of non-linear elliptic equations}
give the systematic definitions and theorem statements. T.9.6.17 and
T.10.2.4 prove the pullback-preservation results. N.5.3 and N.6.5 contain the
proposed analytic and derived regularity arguments
\cite[P.6--P.7, T.9.6.17, T.10.2.4, N.5.3, and N.6.5]
  {Pardon_Logarithmic_Derived_Moduli}. The manuscript is unfinished, and that
status should accompany every citation to its regularity theorems.

Pardon's
\emph{An algebraic approach to virtual fundamental cycles on moduli spaces of
pseudo-holomorphic curves} develops implicit atlases and the algebraic VFC
package. Its introduction is sufficient for the comparison made here
\cite[Sections~1.1--1.3]{Pardon_Virtual_Fundamental_Cycles}. It should not be
used as evidence that derived representability itself produces a virtual
class.

For the earlier derived-smooth viewpoint on bordism and intersections, see
Spivak's \emph{Derived smooth manifolds}
\cite{Spivak_Derived_Smooth_Manifolds}. A precise comparison among derived
bordism, implicit-atlas virtual classes, and other virtual formalisms lies
beyond this seminar.
